\documentclass[a4paper]{article}
\usepackage[fontset=fandol]{ctex}
\usepackage[margin=1.65cm]{geometry}
\usepackage{amsmath,amssymb}
\usepackage{graphicx}
\usepackage{hyperref}
\usepackage{longtable,booktabs}
\usepackage{xcolor}
\usepackage{enumitem}
\usepackage{listings}
\lstset{basicstyle=\ttfamily\small,breaklines=true,columns=fullflexible,keepspaces=true}
\hypersetup{colorlinks=true,linkcolor=blue,urlcolor=blue}
\setlist[itemize]{leftmargin=1.25em,itemsep=0.18em}
\setlength{\parskip}{0.45em}
\setlength{\parindent}{2em}
\sloppy
\title{第 12 讲：Evaluation 方法论}
\author{CS336 Spring 2026 public videos and official course materials}
\date{}
\begin{document}
\maketitle
\tableofcontents
\newpage
\section{本章导读}
本章讨论 \textbf{Evaluation 方法论}。CS336 的第一层目标不是把现成大模型当作黑箱使用，而是把 language model 从输入文本、训练目标、模型结构、数据、系统和评估这些部件重新拆开。只有拆开之后，读者才能判断一个设计选择究竟改变了什么。
设想两个 checkpoint 的 validation loss 很接近，但一个在 coding benchmark 上好、另一个更安全。Evaluation（评估）不是给模型打一个总分，而是建立一套可复现、抗污染、能服务决策的测量系统。
本章对应 CS336 Spring 2026 第 12 个公开视频，并结合课程页提供的可解析资料。视频和课程页给出事实边界；正文把这些材料改写为可自学的教材章节。
阅读本章时，先理解问题为什么存在，再看定义、公式和实现形态，随后通过例题和练习检查自己是否能独立复现这条 reasoning chain。
\begin{figure}[h]
\centering
\includegraphics[width=0.58\linewidth]{source_anchor.jpg}
\caption{第 12 讲的课程图像锚点。}
\end{figure}
\subsection{本章知识路线}
\begin{itemize}
\item 本章首先说明 Evaluation 方法论 为什么是 language modeling pipeline 中不可跳过的一环。
\item 其次定义本章反复使用的术语，并说明这些术语对应的计算对象或系统对象。
\item 随后用公式和伪代码表达主要机制，让概念能够落到可计算的变量上。
\item 最后给出例题、风险讨论和练习，便于读者检查自己是否真正掌握了机制。
\end{itemize}
\subsection{结构图}
下面的图用于帮助读者定位本章的核心结构。先看概念之间的依赖，再回到正文中的公式、伪代码和例题。
\begin{figure}[h]
\centering
\includegraphics[width=0.72\linewidth]{figures/source_grounded_concept_map.png}
\caption{第 12 讲概念图：本章问题、术语和机制的关系。}
\end{figure}
\begin{figure}[h]
\centering
\includegraphics[width=0.72\linewidth]{figures/source_grounded_checklist.png}
\caption{第 12 讲学习路径：从问题定义进入公式、实现、实验和练习。}
\end{figure}
\section{本章主线}
本章可以沿着三条线索阅读：课程为什么把这个主题放在这里，它改变了哪些变量，以及怎样用小实验检查这些变量。
\subsection{本章要解决的核心问题}
本讲讲 evaluation：训练出模型后，必须用 loss、perplexity、benchmark、human preference、安全测试和部署指标分别回答不同问题。
这一点对应的学习动作是：找出对象，写出变量，说明变量如何被测量。若输入规模、硬件、数据分布或评估协议改变，结论也必须重新检查。
因此，Evaluation 方法论 不应被读成若干名词的集合，而应被读成一组可检验的机制。
\subsection{从机制到资源账本}
benchmark 不是一个数字，而是一套 prompt、scoring rule、aggregation、decoding setting 和 contamination control。测试集污染会让分数失去泛化意义。
这一点对应的学习动作是：找出对象，写出变量，说明变量如何被测量。若输入规模、硬件、数据分布或评估协议改变，结论也必须重新检查。
因此，Evaluation 方法论 不应被读成若干名词的集合，而应被读成一组可检验的机制。
\subsection{从课堂讲解到可复现实验}
评估要服务于决策：选 checkpoint、选 data mixture、判断 post-training 是否改善、判断 serving tradeoff 是否可接受。
这一点对应的学习动作是：找出对象，写出变量，说明变量如何被测量。若输入规模、硬件、数据分布或评估协议改变，结论也必须重新检查。
因此，Evaluation 方法论 不应被读成若干名词的集合，而应被读成一组可检验的机制。
\section{术语、符号与问题边界}
本章的术语横跨模型、数据和系统。读者需要先知道每个词指向的对象：有些词描述数学目标，有些词描述张量形状，有些词描述硬件瓶颈，还有一些词描述评估协议。术语边界不清，后面的公式和实现就会变成空洞的缩写。
\begin{longtable}{p{0.08\linewidth}p{0.42\linewidth}p{0.42\linewidth}}
\toprule
\textbf{\#} & \textbf{术语} & \textbf{阅读要求}\\
\midrule
1 & evaluation（评估） & 需要能说明它指向的对象、参与的公式或代码位置，以及一个典型失败情形。\\
2 & perplexity（困惑度） & 需要能说明它指向的对象、参与的公式或代码位置，以及一个典型失败情形。\\
3 & benchmark & 需要能说明它指向的对象、参与的公式或代码位置，以及一个典型失败情形。\\
4 & contamination（污染） & 需要能说明它指向的对象、参与的公式或代码位置，以及一个典型失败情形。\\
5 & calibration（校准） & 需要能说明它指向的对象、参与的公式或代码位置，以及一个典型失败情形。\\
6 & human evaluation（人工评估） & 需要能说明它指向的对象、参与的公式或代码位置，以及一个典型失败情形。\\
\bottomrule
\end{longtable}
这些术语之间有依赖关系。例如 tokenization 改变 sequence length，sequence length 又改变 attention FLOPs、KV cache 和 evaluation token budget；parallelism 改变显存占用，也改变通信瓶颈和故障恢复方式。
\subsection{关键词}
下面的关键词在课程材料中反复出现。它们不代替定义，只提示读者本章哪些对象需要在后文持续跟踪。
\begin{itemize}
\item perplexity, evaluation, language, tasks, function, people, data, prompts, benchmarks, need, e.g., response, evaluating, agent
\end{itemize}
\section{评估在 LM 生命周期中的位置}
现在进入 评估在 LM 生命周期中的位置。本节关心的不是名称本身，而是它把 Evaluation 方法论 中的哪些对象变成可计算、可测量或可优化的量。
\subsection{为什么要讨论 评估在 LM 生命周期中的位置}
视频指出课程至此已覆盖 architecture、optimizer、training loop、systems、scaling 和 inference，剩下的关键是如何判断模型好坏。
评估连接训练目标、数据选择、模型比较和部署决策。
好的评估应区分 intrinsic loss、task benchmark、用户偏好、安全性和系统指标。
评估在 LM 生命周期中的位置 的作用是把 Evaluation 方法论 中的一个抽象问题变成可计算对象：哪些变量进入公式，哪些状态进入代码，哪些条件会让结果失效。
\subsection{评估在 LM 生命周期中的位置 的机制}
视频指出课程至此已覆盖 architecture、optimizer、training loop、systems、scaling 和 inference，剩下的关键是如何判断模型好坏。 这给出本节的对象和边界。读者需要把其中的名词对应到张量、数据记录、训练状态或系统状态，而不是停留在缩写层面。
评估连接训练目标、数据选择、模型比较和部署决策。 因而同一个方法在小规模示例、单机实验和集群训练中的意义并不相同。比较方法时，问题规模、硬件条件、数据分布和评估口径必须同时给出。
好的评估应区分 intrinsic loss、task benchmark、用户偏好、安全性和系统指标。 这使本节具有可检验性：如果一个结论不能在公式、伪代码、profile trace、benchmark 或 ablation 中表现出来，它就还只是直觉。
因此，评估在 LM 生命周期中的位置 应当被理解为一组可以实现和测试的假设。CS336 反复强调 from scratch，正是因为只有能被复现的机制，才可能被稳定改进。
\subsection{评估在 LM 生命周期中的位置 的数学表达}
\[\operatorname{PPL}=\exp\left(-\frac{1}{T}\sum_{t=1}^{T}\log p_\theta(x_t\mid x_{<t})\right)\]
符号说明：PPL 是 token-level cross entropy 的指数形式；越低通常表示语言建模越好。
这个关系式给出了本节最小的数学骨架。读者应检查每个变量的单位、取值范围和可观测性；如果某个量只能间接估计，后续实验就必须说明 proxy 是什么。
例如，validation loss 常被用作模型质量的 proxy，tokens/s 用作吞吐 proxy，Jaccard 或 MinHash 用作近重复 proxy，reward model score 用作偏好 proxy。proxy 不是目标本身；它只在评估协议清楚时才可引用。
\subsection{从 评估在 LM 生命周期中的位置 到程序}
本节机制可以写成下面的伪代码。它保留课程材料中的计算顺序和状态变化，但省略了工程代码中的错误处理、并行细节和框架样板。
\begin{lstlisting}[language=Python]
evaluate_checkpoints = [loss, perplexity, benchmark_suite, safety_suite]
compare_with_confidence_intervals(evaluate_checkpoints)
\end{lstlisting}
读这段伪代码时，重点不是语法，而是状态如何变化：输入是什么，中间变量是什么，主要计算或数据移动发生在哪里，哪些边界条件会破坏结果。
\subsection{边界条件与常见误解}
\begin{itemize}
\item perplexity 不直接等于 instruction-following 能力。
\item 不同 tokenizer 下的 perplexity 不可直接比较。
\item 不要把小实验中的局部结论自动外推到 frontier-scale；scale、数据分布和硬件拓扑都会改变结论。
\item 不要把 benchmark 或 profile 的单个数字当作机制解释；必须同时记录实验设置、输入形状、版本和评估口径。
\end{itemize}
\subsection{与本章其他概念的关系}
评估在 LM 生命周期中的位置 与本章其他部分的关系是互补的：术语表给出对象名称，公式给出最小数学结构，伪代码给出执行顺序，例题则把这些内容压缩成一个可检查的计算。
\subsection{怎样检验 评估在 LM 生命周期中的位置}
检验 评估在 LM 生命周期中的位置 的第一步是确定测量对象。这里的测量对象通常不是一个抽象名词，而是与 evaluation、perplexity、loss、benchmark、checkpoint 有关的张量、样本、请求、奖励、通信操作或评估记录。如果对象没有被具体化，后面的实验就只能得到描述性观察，不能得到可复现结论。
第二步是构造最小输入。最小输入应该足够小，使读者能手算或打印关键中间量；同时又要保留本节机制。对于 评估在 LM 生命周期中的位置，最小输入不追求逼真规模，而追求暴露因果关系：改变一个变量时，哪些输出、成本或错误会随之改变。
第三步是写出资源账本。账本可以是 FLOPs、bytes moved、显存峰值、通信量、token 数、样本数、reward 分布或 benchmark 分数。符号说明：PPL 是 token-level cross entropy 的指数形式；越低通常表示语言建模越好。 这类说明应被转化为单位明确的数量关系；如果数量只能间接观测，就必须说明使用了什么 proxy。
第四步是设置对照。一个好对照只改变一个因素，例如 batch size、sequence length、vocabulary size、attention heads、filter threshold、learning rate、decoding 参数或 verifier。如果多个因素同时改变，实验最多说明系统表现不同，不能说明是哪一个机制在起作用。
第五步是观察失败条件。教材中的成功路径往往最容易记住，但工程中真正决定方法边界的是失败案例。本节至少要记住这些限制：perplexity 不直接等于 instruction-following 能力；不同 tokenizer 下的 perplexity 不可直接比较。复习时可以把这些限制改写成反例测试。
最后，把实验结论放回 Evaluation 方法论 的整体结构中。评估在 LM 生命周期中的位置 不是孤立技巧：它会影响训练目标、数据选择、系统成本或评估解释中的至少一项。能说清楚这种影响，才说明读者已经从名词记忆进入机制理解。
\subsection{常见混淆：评估在 LM 生命周期中的位置}
evaluation（评估）把 loss、benchmark、人类偏好和部署指标区分开。在“评估在 LM 生命周期中的位置”中，这个概念用于解释 Evaluation 方法论 的一个具体机制。
perplexity 在本节中应被看作一个技术对象，而不是一个孤立名词。需要判断它指向模型、数据、系统还是评估，并说明它在“评估在 LM 生命周期中的位置”中改变了哪个变量或约束。
loss 在本节中应被看作一个技术对象，而不是一个孤立名词。需要判断它指向模型、数据、系统还是评估，并说明它在“评估在 LM 生命周期中的位置”中改变了哪个变量或约束。
benchmark 在本节中应被看作一个技术对象，而不是一个孤立名词。需要判断它指向模型、数据、系统还是评估，并说明它在“评估在 LM 生命周期中的位置”中改变了哪个变量或约束。
checkpoint 在本节中应被看作一个技术对象，而不是一个孤立名词。需要判断它指向模型、数据、系统还是评估，并说明它在“评估在 LM 生命周期中的位置”中改变了哪个变量或约束。
这些概念的共同点是：只有进入公式、代码或评估协议之后才有明确含义。单独背诵名称不会帮助判断一个训练 run、推理系统或数据 pipeline 是否正确。
\subsection{评估在 LM 生命周期中的位置 的小结}
\begin{itemize}
\item 评估在 LM 生命周期中的位置 的任何性能结论都必须同时报告问题规模和评估口径，否则不可比较。
\item 如果一个方法声称降低主要成本，它通常会把成本转移到另一个变量，例如显存、通信、数据质量、训练稳定性或评估复杂度。
\item 公式和伪代码提供推理骨架；真实系统仍需要 profiler、ablation 和 held-out evaluation。
\end{itemize}
\section{Benchmark 设计与污染}
现在进入 Benchmark 设计与污染。本节关心的不是名称本身，而是它把 Evaluation 方法论 中的哪些对象变成可计算、可测量或可优化的量。
\subsection{为什么要讨论 Benchmark 设计与污染}
benchmark 要明确定义任务、数据、prompt、few-shot examples、scoring rule 和 aggregation。
contamination 是评估大模型时的核心风险：测试集可能进入 pretraining 或 post-training 数据。
评估应记录 exact prompt、decoding parameters、model checkpoint 和版本，保证可复现。
Benchmark 设计与污染 的作用是把 Evaluation 方法论 中的一个抽象问题变成可计算对象：哪些变量进入公式，哪些状态进入代码，哪些条件会让结果失效。
\subsection{Benchmark 设计与污染 的机制}
benchmark 要明确定义任务、数据、prompt、few-shot examples、scoring rule 和 aggregation。 这给出本节的对象和边界。读者需要把其中的名词对应到张量、数据记录、训练状态或系统状态，而不是停留在缩写层面。
contamination 是评估大模型时的核心风险：测试集可能进入 pretraining 或 post-training 数据。 因而同一个方法在小规模示例、单机实验和集群训练中的意义并不相同。比较方法时，问题规模、硬件条件、数据分布和评估口径必须同时给出。
评估应记录 exact prompt、decoding parameters、model checkpoint 和版本，保证可复现。 这使本节具有可检验性：如果一个结论不能在公式、伪代码、profile trace、benchmark 或 ablation 中表现出来，它就还只是直觉。
因此，Benchmark 设计与污染 应当被理解为一组可以实现和测试的假设。CS336 反复强调 from scratch，正是因为只有能被复现的机制，才可能被稳定改进。
\subsection{Benchmark 设计与污染 的数学表达}
\[\hat s=\frac{1}{n}\sum_{i=1}^{n}\mathbf{1}\{f_\theta(q_i)=a_i\}\]
符号说明：\(\hat s\) 是平均准确率；\(q_i,a_i\) 是第 \(i\) 个样本的问题和答案。
这个关系式给出了本节最小的数学骨架。读者应检查每个变量的单位、取值范围和可观测性；如果某个量只能间接估计，后续实验就必须说明 proxy 是什么。
例如，validation loss 常被用作模型质量的 proxy，tokens/s 用作吞吐 proxy，Jaccard 或 MinHash 用作近重复 proxy，reward model score 用作偏好 proxy。proxy 不是目标本身；它只在评估协议清楚时才可引用。
\subsection{从 Benchmark 设计与污染 到程序}
本节机制可以写成下面的伪代码。它保留课程材料中的计算顺序和状态变化，但省略了工程代码中的错误处理、并行细节和框架样板。
\begin{lstlisting}[language=Python]
for item in benchmark:
    prompt = render_prompt(template, item)
    pred = model.generate(prompt, decoding)
    score += grader(pred, item.answer)
\end{lstlisting}
读这段伪代码时，重点不是语法，而是状态如何变化：输入是什么，中间变量是什么，主要计算或数据移动发生在哪里，哪些边界条件会破坏结果。
\subsection{边界条件与常见误解}
\begin{itemize}
\item 自动 grader 可能奖励格式匹配而非真实能力。
\item 测试集泄漏会让 score 失去泛化意义。
\item 不要把小实验中的局部结论自动外推到 frontier-scale；scale、数据分布和硬件拓扑都会改变结论。
\item 不要把 benchmark 或 profile 的单个数字当作机制解释；必须同时记录实验设置、输入形状、版本和评估口径。
\end{itemize}
\subsection{与本章其他概念的关系}
Benchmark 设计与污染 与本章其他部分的关系是互补的：术语表给出对象名称，公式给出最小数学结构，伪代码给出执行顺序，例题则把这些内容压缩成一个可检查的计算。
\subsection{怎样检验 Benchmark 设计与污染}
检验 Benchmark 设计与污染 的第一步是确定测量对象。这里的测量对象通常不是一个抽象名词，而是与 benchmark、prompt、grader、contamination、reproducible 有关的张量、样本、请求、奖励、通信操作或评估记录。如果对象没有被具体化，后面的实验就只能得到描述性观察，不能得到可复现结论。
第二步是构造最小输入。最小输入应该足够小，使读者能手算或打印关键中间量；同时又要保留本节机制。对于 Benchmark 设计与污染，最小输入不追求逼真规模，而追求暴露因果关系：改变一个变量时，哪些输出、成本或错误会随之改变。
第三步是写出资源账本。账本可以是 FLOPs、bytes moved、显存峰值、通信量、token 数、样本数、reward 分布或 benchmark 分数。符号说明：\textbackslash{}(\textbackslash{}hat s\textbackslash{}) 是平均准确率；\textbackslash{}(q\_i,a\_i\textbackslash{}) 是第 \textbackslash{}(i\textbackslash{}) 个样本的问题和答案。 这类说明应被转化为单位明确的数量关系；如果数量只能间接观测，就必须说明使用了什么 proxy。
第四步是设置对照。一个好对照只改变一个因素，例如 batch size、sequence length、vocabulary size、attention heads、filter threshold、learning rate、decoding 参数或 verifier。如果多个因素同时改变，实验最多说明系统表现不同，不能说明是哪一个机制在起作用。
第五步是观察失败条件。教材中的成功路径往往最容易记住，但工程中真正决定方法边界的是失败案例。本节至少要记住这些限制：自动 grader 可能奖励格式匹配而非真实能力；测试集泄漏会让 score 失去泛化意义。复习时可以把这些限制改写成反例测试。
最后，把实验结论放回 Evaluation 方法论 的整体结构中。Benchmark 设计与污染 不是孤立技巧：它会影响训练目标、数据选择、系统成本或评估解释中的至少一项。能说清楚这种影响，才说明读者已经从名词记忆进入机制理解。
\subsection{常见混淆：Benchmark 设计与污染}
benchmark 在本节中应被看作一个技术对象，而不是一个孤立名词。需要判断它指向模型、数据、系统还是评估，并说明它在“Benchmark 设计与污染”中改变了哪个变量或约束。
prompt 在本节中应被看作一个技术对象，而不是一个孤立名词。需要判断它指向模型、数据、系统还是评估，并说明它在“Benchmark 设计与污染”中改变了哪个变量或约束。
grader 在本节中应被看作一个技术对象，而不是一个孤立名词。需要判断它指向模型、数据、系统还是评估，并说明它在“Benchmark 设计与污染”中改变了哪个变量或约束。
contamination 在本节中应被看作一个技术对象，而不是一个孤立名词。需要判断它指向模型、数据、系统还是评估，并说明它在“Benchmark 设计与污染”中改变了哪个变量或约束。
reproducible 在本节中应被看作一个技术对象，而不是一个孤立名词。需要判断它指向模型、数据、系统还是评估，并说明它在“Benchmark 设计与污染”中改变了哪个变量或约束。
这些概念的共同点是：只有进入公式、代码或评估协议之后才有明确含义。单独背诵名称不会帮助判断一个训练 run、推理系统或数据 pipeline 是否正确。
\subsection{Benchmark 设计与污染 的小结}
\begin{itemize}
\item Benchmark 设计与污染 的任何性能结论都必须同时报告问题规模和评估口径，否则不可比较。
\item 如果一个方法声称降低主要成本，它通常会把成本转移到另一个变量，例如显存、通信、数据质量、训练稳定性或评估复杂度。
\item 公式和伪代码提供推理骨架；真实系统仍需要 profiler、ablation 和 held-out evaluation。
\end{itemize}
\section{从离线评估到 deployment 指标}
现在进入 从离线评估到 deployment 指标。本节关心的不是名称本身，而是它把 Evaluation 方法论 中的哪些对象变成可计算、可测量或可优化的量。
\subsection{为什么要讨论 从离线评估到 deployment 指标}
真实模型上线需要结合 latency、cost、toxicity、helpfulness、truthfulness、robustness 和 refusal behavior。
人工评估和 pairwise preference 能捕捉开放式输出质量，但成本高、方差大、标准易漂移。
课程后续 post-training 会把 evaluation 与 reward modeling/RL 联系起来。
从离线评估到 deployment 指标 的作用是把 Evaluation 方法论 中的一个抽象问题变成可计算对象：哪些变量进入公式，哪些状态进入代码，哪些条件会让结果失效。
\subsection{从离线评估到 deployment 指标 的机制}
真实模型上线需要结合 latency、cost、toxicity、helpfulness、truthfulness、robustness 和 refusal behavior。 这给出本节的对象和边界。读者需要把其中的名词对应到张量、数据记录、训练状态或系统状态，而不是停留在缩写层面。
人工评估和 pairwise preference 能捕捉开放式输出质量，但成本高、方差大、标准易漂移。 因而同一个方法在小规模示例、单机实验和集群训练中的意义并不相同。比较方法时，问题规模、硬件条件、数据分布和评估口径必须同时给出。
课程后续 post-training 会把 evaluation 与 reward modeling/RL 联系起来。 这使本节具有可检验性：如果一个结论不能在公式、伪代码、profile trace、benchmark 或 ablation 中表现出来，它就还只是直觉。
因此，从离线评估到 deployment 指标 应当被理解为一组可以实现和测试的假设。CS336 反复强调 from scratch，正是因为只有能被复现的机制，才可能被稳定改进。
\subsection{从离线评估到 deployment 指标 的数学表达}
\[\Pr(A\succ B)=\sigma(r_\phi(x,y_A)-r_\phi(x,y_B))\]
符号说明：preference model 用 reward 差预测人类更偏好哪个回答；\(\sigma\) 是 sigmoid。
这个关系式给出了本节最小的数学骨架。读者应检查每个变量的单位、取值范围和可观测性；如果某个量只能间接估计，后续实验就必须说明 proxy 是什么。
例如，validation loss 常被用作模型质量的 proxy，tokens/s 用作吞吐 proxy，Jaccard 或 MinHash 用作近重复 proxy，reward model score 用作偏好 proxy。proxy 不是目标本身；它只在评估协议清楚时才可引用。
\subsection{从 从离线评估到 deployment 指标 到程序}
本节机制可以写成下面的伪代码。它保留课程材料中的计算顺序和状态变化，但省略了工程代码中的错误处理、并行细节和框架样板。
\begin{lstlisting}[language=Python]
collect_pairwise_preferences(outputs)
fit_reward_model(preferences)
validate_against_heldout_human_judgments()
\end{lstlisting}
读这段伪代码时，重点不是语法，而是状态如何变化：输入是什么，中间变量是什么，主要计算或数据移动发生在哪里，哪些边界条件会破坏结果。
\subsection{边界条件与常见误解}
\begin{itemize}
\item 人工偏好不是绝对真理，会受标注指南和人群影响。
\item reward model 可能被优化过程 exploited。
\item 不要把小实验中的局部结论自动外推到 frontier-scale；scale、数据分布和硬件拓扑都会改变结论。
\item 不要把 benchmark 或 profile 的单个数字当作机制解释；必须同时记录实验设置、输入形状、版本和评估口径。
\end{itemize}
\subsection{与本章其他概念的关系}
从离线评估到 deployment 指标 与本章其他部分的关系是互补的：术语表给出对象名称，公式给出最小数学结构，伪代码给出执行顺序，例题则把这些内容压缩成一个可检查的计算。
\subsection{怎样检验 从离线评估到 deployment 指标}
检验 从离线评估到 deployment 指标 的第一步是确定测量对象。这里的测量对象通常不是一个抽象名词，而是与 human evaluation、preference、reward、deployment、safety 有关的张量、样本、请求、奖励、通信操作或评估记录。如果对象没有被具体化，后面的实验就只能得到描述性观察，不能得到可复现结论。
第二步是构造最小输入。最小输入应该足够小，使读者能手算或打印关键中间量；同时又要保留本节机制。对于 从离线评估到 deployment 指标，最小输入不追求逼真规模，而追求暴露因果关系：改变一个变量时，哪些输出、成本或错误会随之改变。
第三步是写出资源账本。账本可以是 FLOPs、bytes moved、显存峰值、通信量、token 数、样本数、reward 分布或 benchmark 分数。符号说明：preference model 用 reward 差预测人类更偏好哪个回答；\textbackslash{}(\textbackslash{}sigma\textbackslash{}) 是 sigmoid。 这类说明应被转化为单位明确的数量关系；如果数量只能间接观测，就必须说明使用了什么 proxy。
第四步是设置对照。一个好对照只改变一个因素，例如 batch size、sequence length、vocabulary size、attention heads、filter threshold、learning rate、decoding 参数或 verifier。如果多个因素同时改变，实验最多说明系统表现不同，不能说明是哪一个机制在起作用。
第五步是观察失败条件。教材中的成功路径往往最容易记住，但工程中真正决定方法边界的是失败案例。本节至少要记住这些限制：人工偏好不是绝对真理，会受标注指南和人群影响；reward model 可能被优化过程 exploited。复习时可以把这些限制改写成反例测试。
最后，把实验结论放回 Evaluation 方法论 的整体结构中。从离线评估到 deployment 指标 不是孤立技巧：它会影响训练目标、数据选择、系统成本或评估解释中的至少一项。能说清楚这种影响，才说明读者已经从名词记忆进入机制理解。
\subsection{常见混淆：从离线评估到 deployment 指标}
evaluation（评估）把 loss、benchmark、人类偏好和部署指标区分开。在“从离线评估到 deployment 指标”中，这个概念用于解释 Evaluation 方法论 的一个具体机制。
preference 在本节中应被看作一个技术对象，而不是一个孤立名词。需要判断它指向模型、数据、系统还是评估，并说明它在“从离线评估到 deployment 指标”中改变了哪个变量或约束。
reward 在本节中应被看作一个技术对象，而不是一个孤立名词。需要判断它指向模型、数据、系统还是评估，并说明它在“从离线评估到 deployment 指标”中改变了哪个变量或约束。
deployment 在本节中应被看作一个技术对象，而不是一个孤立名词。需要判断它指向模型、数据、系统还是评估，并说明它在“从离线评估到 deployment 指标”中改变了哪个变量或约束。
safety 在本节中应被看作一个技术对象，而不是一个孤立名词。需要判断它指向模型、数据、系统还是评估，并说明它在“从离线评估到 deployment 指标”中改变了哪个变量或约束。
这些概念的共同点是：只有进入公式、代码或评估协议之后才有明确含义。单独背诵名称不会帮助判断一个训练 run、推理系统或数据 pipeline 是否正确。
\subsection{从离线评估到 deployment 指标 的小结}
\begin{itemize}
\item 从离线评估到 deployment 指标 的任何性能结论都必须同时报告问题规模和评估口径，否则不可比较。
\item 如果一个方法声称降低主要成本，它通常会把成本转移到另一个变量，例如显存、通信、数据质量、训练稳定性或评估复杂度。
\item 公式和伪代码提供推理骨架；真实系统仍需要 profiler、ablation 和 held-out evaluation。
\end{itemize}
\section{综合讨论：从局部机制到完整系统}
本章的主题是 Evaluation 方法论，但它不应被读成几个相邻知识点的拼接。评估在 LM 生命周期中的位置、Benchmark 设计与污染、从离线评估到 deployment 指标 共同回答同一个问题：语言模型系统中哪些对象可以被定义，哪些成本可以被估算，哪些结论必须通过实验或评估来确认。
这些对象包括 evaluation（评估）、perplexity（困惑度）、benchmark、contamination（污染）、calibration（校准）、human evaluation（人工评估）。读者可以把它们看成一组接口：有的接口连接文本和 token，有的连接模型结构和张量计算，有的连接数据样本和训练目标，有的连接离线评估和在线服务。接口一旦改变，后面的成本、错误类型和优化空间也会随之改变。
\subsection{评估在 LM 生命周期中的位置 在系统中的位置}
评估在 LM 生命周期中的位置 首先提供一个局部解释：视频指出课程至此已覆盖 architecture、optimizer、training loop、systems、scaling 和 inference，剩下的关键是如何判断模型好坏。 这类解释的价值在于把直觉压缩成可操作对象。一旦对象明确，读者就可以追问它的输入、输出、中间状态和失败条件，而不是只记住方法名。
其次，评估在 LM 生命周期中的位置 会改变至少一个资源或评估变量。它可能改变 sequence length、activation size、memory bandwidth、communication volume、data mixture、reward distribution 或 benchmark protocol。这些变量在小例子里可以手算，在真实训练和服务系统里则需要 profiling、logging 和 held-out evaluation。
最后，评估在 LM 生命周期中的位置 不能脱离边界条件理解。一个关键 caveat 是：perplexity 不直接等于 instruction-following 能力。。把 caveat 写出来不是为了削弱结论，而是为了说明结论在哪些输入、硬件、数据和评估条件下才成立。
\subsection{Benchmark 设计与污染 在系统中的位置}
Benchmark 设计与污染 首先提供一个局部解释：benchmark 要明确定义任务、数据、prompt、few-shot examples、scoring rule 和 aggregation。 这类解释的价值在于把直觉压缩成可操作对象。一旦对象明确，读者就可以追问它的输入、输出、中间状态和失败条件，而不是只记住方法名。
其次，Benchmark 设计与污染 会改变至少一个资源或评估变量。它可能改变 sequence length、activation size、memory bandwidth、communication volume、data mixture、reward distribution 或 benchmark protocol。这些变量在小例子里可以手算，在真实训练和服务系统里则需要 profiling、logging 和 held-out evaluation。
最后，Benchmark 设计与污染 不能脱离边界条件理解。一个关键 caveat 是：自动 grader 可能奖励格式匹配而非真实能力。。把 caveat 写出来不是为了削弱结论，而是为了说明结论在哪些输入、硬件、数据和评估条件下才成立。
\subsection{从离线评估到 deployment 指标 在系统中的位置}
从离线评估到 deployment 指标 首先提供一个局部解释：真实模型上线需要结合 latency、cost、toxicity、helpfulness、truthfulness、robustness 和 refusal behavior。 这类解释的价值在于把直觉压缩成可操作对象。一旦对象明确，读者就可以追问它的输入、输出、中间状态和失败条件，而不是只记住方法名。
其次，从离线评估到 deployment 指标 会改变至少一个资源或评估变量。它可能改变 sequence length、activation size、memory bandwidth、communication volume、data mixture、reward distribution 或 benchmark protocol。这些变量在小例子里可以手算，在真实训练和服务系统里则需要 profiling、logging 和 held-out evaluation。
最后，从离线评估到 deployment 指标 不能脱离边界条件理解。一个关键 caveat 是：人工偏好不是绝对真理，会受标注指南和人群影响。。把 caveat 写出来不是为了削弱结论，而是为了说明结论在哪些输入、硬件、数据和评估条件下才成立。
\subsection{把本章知识用于读论文和复现实验}
读相关论文时，可以把本章变成一个检查表。第一，论文是否清楚定义了输入规模、模型规模、数据来源和评估协议；第二，论文声称的改进落在哪个变量上；第三，作者是否报告了会让结论失效的设置；第四，实验是否能分离模型结构、数据质量、优化设置和系统实现的影响。
复现实验时，也应避免直接追求大规模。先用最小程序复现一个公式、一个伪代码分支或一个 profile 现象，再扩大输入规模。这样做的原因很简单：如果小规模程序无法解释中间量，大规模训练只会把错误隐藏在日志和随机波动里。
把这些检查合在一起，本章给出的不是一个固定答案，而是一种阅读 CS336 的方法：每个模型机制都要落到变量，每个变量都要落到测量，每个测量都要说明协议，每个协议都要承认边界。
因此，本章中的任何一句结论都可以被改写成一个可引用的完整句式：在给定数据、模型、硬件和评估协议下，某个机制改变了某个变量，并通过某个观测指标表现出来。这样的句子比单独背诵术语更长，但它包含了自学、复现和引用时真正需要的信息。
把本章用于自学时，可以按三遍阅读。第一遍只追踪对象和定义，确认每个术语到底指向文本、token、张量、样本、请求还是评估记录。第二遍追踪公式和伪代码，确认每一步改变了哪些状态。第三遍追踪实验和 caveat，确认哪些结论只在特定规模、数据或硬件条件下成立。
把本章用于复习时，则应反过来操作：先从一个失败案例或异常指标出发，再回到相应机制。loss 曲线异常可能指向数据、优化或模型容量；吞吐异常可能指向 kernel、通信或调度；benchmark 异常可能指向 contamination、prompt protocol 或采样设置。能够从异常反推机制，是掌握本章的实用标准。
\section{例题：把概念变成可计算检查}
\subsection{同一 benchmark 的可复现评分}
评估一个 multiple-choice benchmark。
\begin{itemize}
\item 固定 prompt template、few-shot examples、答案解析方式和 decoding 参数。
\item 记录 tokenizer、checkpoint、commit hash 和数据版本。
\item 报告平均分时给出方差或置信区间，并检查 contamination。
\end{itemize}
\textbf{结论。} benchmark 分数只有在评估协议完整时才可引用。
这个例题的目的不是替代完整工程实现，而是给出一种最小推理形式：把问题转成变量，写出近似公式或伪代码，再说明近似何时失效。
\section{实验设计：从小例子到可复现结论}
学习 Evaluation 方法论 不能停在概念层。一个教材化结论至少需要四件事：可构造的输入、可观测的中间量、可比较的指标，以及清楚的失败条件。下面的实验设计不是要求训练大模型，而是要求读者用小规模程序复现本章机制。
\subsection{最小输入}
先构造只触发本章核心机制的小输入。例如 tokenizer 章节可以用几个包含 rare words 和 Unicode 字符的字符串；GPU 章节可以用一个 elementwise kernel 和一个 GEMM；data 章节可以用十几篇近重复文档。小输入的价值在于它让错误可见。
\subsection{资源账本}
对同一输入写下 FLOPs、bytes moved、显存峰值、通信量或样本数量的估算。估算不需要一开始就精确，但必须说明单位和近似假设。随后用 profiler、benchmark、validation loss 或人工检查去验证估算是否同量级。
\subsection{单变量消融}
一次只改变一个因素，例如 vocabulary size、head 数、batch size、filter threshold、reward scale、decoding temperature 或 GPU 数量。若多个变量同时变化，实验只能说明系统变了，不能说明机制是什么。
\subsection{失败条件}
最后要刻意触发失败：极端 context length、错误 dtype、过小 batch、污染 benchmark、低质量数据或不可靠 verifier。失败条件常常比成功曲线更能说明一个方法的边界。
\subsection{记录与报告}
实验记录应能让另一个读者复现同一结论。最少要写清楚输入构造、代码版本、随机种子、硬件、batch shape、数据版本、评价指标和异常处理方式。报告结论时不要只写“更快”或“更好”，而要说明快在哪里、好在哪个指标上，以及是否牺牲了内存、稳定性、数据质量或可解释性。
常见报告错误是把局部现象写成普遍规律。例如一次小模型实验里的 loss 改善，不等于更大模型或不同数据分布上也会改善；一次 kernel benchmark 的吞吐提升，不等于端到端 serving latency 一定降低。教材中的实验报告应始终把结论限定在可复现条件内。
\section{课程资料与回看路径}
教材正文已经把主要概念、公式和实现逻辑组织成连续叙述；本节只提供回到课程材料的阅读路径。读者复习时可以先完成前文练习，再按下面的时间段或资料组核对细节。
\subsection{视频回看路径}
\begin{longtable}{p{0.20\linewidth}p{0.46\linewidth}p{0.26\linewidth}}
\toprule
\textbf{时间} & \textbf{关键词} & \textbf{复习任务}\\
\midrule
00:00:05-00:06:05 & good, evaluation, going, people, data, maybe, different & 回看定义、例子、推导或 caveat 在该段中的位置。\\
00:18:38-00:25:01 & then, language, think, it's, question, benchmark, question. & 回看定义、例子、推导或 caveat 在该段中的位置。\\
00:37:31-00:44:21 & alpacaeval, evaluate, there's, don't, metric, chatbot, arena & 回看定义、例子、推导或 caveat 在该段中的位置。\\
00:51:41-00:58:20 & agent, it's, language, then, reasoning, this., context & 回看定义、例子、推导或 caveat 在该段中的位置。\\
01:11:54-01:18:10 & there's, evaluation, think, very, benchmarks, always, cases & 回看定义、例子、推导或 caveat 在该段中的位置。\\
\bottomrule
\end{longtable}
这些时间段用于定位视频中的讲解顺序。严肃引用时，应回到视频片段确认讲者表述、板书或幻灯片内容。
回看视频时，不必逐字抄录字幕。更有效的方法是把每个时间段判定为定义、例子、推导、代码解释或 caveat，再把它放回本章对应小节。
\subsection{课程资料阅读路径}
\begin{longtable}{p{0.14\linewidth}p{0.48\linewidth}p{0.30\linewidth}}
\toprule
\textbf{资料组} & \textbf{标题/页块} & \textbf{关键词}\\
\midrule
official\_01 & Lecture 12: evaluation & evaluation, good, maybe, training, data, shapes, behavior\\
official\_03 & - For downstream tasks, `response = LM(prompt)` and compute accuracy on `response` & questions, accuracy, language, response, summary, still, benchmarks\\
official\_05 & - Problem: LLM judges favor longer responses, resulted in leaderboard gaming & tasks, evaluate, responses, judge, judges, metric, chatbot\\
official\_08 & Route 3: use fresh evals & evaluating, route, evals, aren, either, perplexity, evaluation\\
official\_10 & - Fixed up SWE-Bench to produce SWE-Bench Verified & function, swe-bench, benchmarks, problems, agent, fixed, produce\\
\bottomrule
\end{longtable}
课程资料通常给出更稳定的标题、公式、图或代码结构；视频则补充动机和口头解释。两者合起来构成本章的事实边界。
阅读资料时，可以把每个资料组拆成三个对象：一个定义，一个公式或伪代码动作，一个边界条件。这样比逐字翻译页面或脚本更容易形成可用知识。
\section{实践说明、历史位置与风险}
Evaluation 方法论 在 CS336 的课程结构中连接基础机制和规模化问题。基础机制解释方法为什么成立；规模化问题检验它在更大模型、更长上下文、更复杂数据或真实 serving workload 下是否仍成立。
历史位置不等于模型年表。更重要的是识别问题如何迁移：早期方法解决小规模建模问题，现代方法常常解决同一问题在大规模数据、GPU 集群、长上下文或人类偏好约束下的新形态。
\begin{itemize}
\item 资料版本限制（version risk）：课程页、公开视频和配套资料可能在学期中更新；引用时应注明访问日期或资料版本。
\item 测量口径限制（measurement risk）：FLOPs、tokens/s、benchmark score、reward 等数字只有在协议完整时可比较。
\item 规模外推限制（scaling risk）：小模型或小 batch 的机制不必然外推到 frontier-scale。
\item 数据分布限制（data risk）：数据来源、过滤和去重策略会改变模型行为，也会改变 evaluation contamination 风险。
\end{itemize}
\section{复习要点}
下面的问题把本章内容压缩为可反复检查的条目。每个条目都应能回到前文的解释、公式、伪代码或例题。
\begin{longtable}{p{0.16\linewidth}p{0.38\linewidth}p{0.38\linewidth}}
\toprule
\textbf{条目} & \textbf{对象} & \textbf{检查动作}\\
\midrule
条目 1 & evaluation（评估） & 定义它；写出一个最小例子；说明一个 failure mode；指出它影响哪个资源或评估账本。\\
条目 2 & perplexity（困惑度） & 定义它；写出一个最小例子；说明一个 failure mode；指出它影响哪个资源或评估账本。\\
条目 3 & benchmark & 定义它；写出一个最小例子；说明一个 failure mode；指出它影响哪个资源或评估账本。\\
条目 4 & contamination（污染） & 定义它；写出一个最小例子；说明一个 failure mode；指出它影响哪个资源或评估账本。\\
条目 5 & calibration（校准） & 定义它；写出一个最小例子；说明一个 failure mode；指出它影响哪个资源或评估账本。\\
条目 6 & human evaluation（人工评估） & 定义它；写出一个最小例子；说明一个 failure mode；指出它影响哪个资源或评估账本。\\
\bottomrule
\end{longtable}
\subsection{综合自测}
\begin{itemize}
\item 我能否不用课程原句，独立解释本章中心问题？
\item 我能否把本章至少一个公式改写成可运行的估算函数？
\item 我能否指出一个看似改进但实际转移成本的方法？
\item 我能否回到视频和课程资料，找到支持本章关键论断的位置？
\item 我能否设计一个最小 ablation，验证本章一个 caveat？
\end{itemize}
\section{总结与延伸}
\subsection{本章小结}
\begin{itemize}
\item Evaluation 方法论 的核心不是记忆术语，而是理解对象、变量和机制之间的关系。
\item 视频给出讲解顺序，课程资料提供可核对的公式、图、代码或页面结构。
\item 复习时应把每个术语连接到至少一个变量、一个实现动作和一个失败模式。
\end{itemize}
\subsection{延伸解释}
本章中的延伸解释用于连接课程材料中的概念，例如说明公式里的 proxy、解释系统 tradeoff，或把多个材料片段串成一个完整推理链。若需要严肃引用，应回到课程视频和配套资料核对原始表述。
\subsection{练习}
\begin{itemize}
\item 定义题：用中英双语解释本章任意五个重要术语，并说明它们属于 compute、memory、data、optimization、evaluation 还是 deployment 账本。
\item 推导题：选择本章一个公式，逐个解释符号、单位和近似假设，并说明哪个变量最难在真实系统中测量。
\item 实现题：把本章伪代码改写成一个最小 Python sanity check，不要求训练大模型，但要输出可检查的中间量。
\item 资料题：从“课程资料与回看路径”中选择一个视频时间段，回到原始视频核对讲者例子，并记录是否存在字幕误差。
\item 评估题：为本章方法设计一个 ablation 或 benchmark，说明控制变量、数据版本、指标和 failure criteria。
\item 边界题：说明本章一个结论在更长 context、更大 batch、更多 GPU 或更复杂数据分布下可能如何失效。
\item 综合题：把本章内容和前一章或后一章连接起来，写出一个端到端训练/推理 pipeline 中的依赖关系。
\item 批判题：指出一个容易被营销材料夸大的说法，并用本章的资源账本或评估协议反驳它。
\end{itemize}
\end{document}
